\section{公的年金の財政}

\begin{fillblank}[年金2 2019 問題1(3)]
  次は公的年金に関する給付水準調整期間の推計の技術的な流れ、および基礎年金の財政の仕
組みの概略についての記述である。以下の空欄を埋めなさい。

\;

（　　A　　）事業の財政において、 （　　B　　）の終了時に給付の支給に支障が生じないようにするため
に必要な積立金を保有しつつ均衡が保たれるように、（　　C　　）部分の給付水準調整期間の終了
年度の見通しが定まる。これにより将来における（　　C　　）部分の水準の見通しが決まる。
将来の（　　C　　）部分の水準を踏まえて、（　　D　　）事業の財政において、 （　　B　　）の終了時に給付の
支給に支障が生じないようにするために必要な積立金を保有しつつ均衡が保たれるように、
（　　E　　）の給付水準調整期間の終了年度の見通しが定まる。これにより、将来の給付水準（所
得代替率）の見通しが作成される。

\;

基礎年金の財政については、毎年度の基礎年金給付に要する費用をその年度の各制度・実施
機関からの（　　F　　）による収入で賄う方式で運営することとされており、各制度・実施機関が
負担する（　　F　　）は、（　　G　　）に応じて按分することとされている。また、公的年金の財政検証
においては、
\begin{itemize}
  \item 毎年度の基礎年金に要する費用とその（　　H　　）
  \item 毎年度の支出を賄うために各公的年金制度が負担すべき（　　F　　）
が推計されている。
\end{itemize}

\end{fillblank}

\begin{answerbox}
  \begin{enumerate}[A: ]
    \item 国民年金
    \item 財政均衡期間
    \item 基礎年金
    \item 厚生年金保険
    \item 報酬比例部分
    \item 基礎年金拠出金
    \item 被保険者の人数
    \item 国庫負担
  \end{enumerate}
\end{answerbox}



\begin{fillblank}[数理人試験 2025 問題1 設問7]
  \noindent \textbf{◯厚生年金保険法}
  \begin{lawstructure}
    \begin{lawarticle}{第八十条}{国庫負担等}
      国庫は、毎年度、厚生年金保険の実施者たる政府が負担する（　　A　　）の額の（　　B　　）に相当する額を負担する。

      \lawparagraph{2}
      国庫は、前項に規定する費用のほか、毎年度、予算の範囲内で、厚生年金保険事業の事務（（　　A　　）の負担に関する事務を含む。次項において同じ。）の執行（実施機関（厚生労働大臣を除く。）によるものを除く。）に要する費用を負担する。

      \lawparagraph{3}
      実施機関（厚生労働大臣を除く。以下この項において同じ。）が納付する（　　A　　）及び実施機関による厚生年金保険事業の事務の執行に要する費用の負担については、この法律に定めるもののほか、（　　C　　）の定めるところによる。
    \end{lawarticle}

    \begin{lawarticle}{第八十四条の五}{拠出金及び政府の負担}
      実施機関は、毎年度、拠出金を納付する。

      \lawparagraph{2}
      次条第一項に規定する（　　D　　）から前項の規定により実施機関が納付する拠出金の合計額及び政府等が負担し、
      又は納付する（　　A　　）保険料相当分（略）の合計額を控除した額については、厚生年金保険の実施者たる政府の負担とする。

      \lawparagraph{3}
      財政の現況及び見通しが作成されるときは、厚生労働大臣は、第一項の規定による実施機関が納付すべき拠出金及び前項の規定による政府の負担について、その将来にわたる予想額を算定するものとする。
    \end{lawarticle}

    \begin{lawarticle}{第八十四条の六}{拠出金の額}
      前条第一項の規定により実施機関が納付する拠出金の額は、当該年度における（　　D　　）に、
      それぞれ次に掲げる率を乗じて得た額の合計額から、当該実施機関が納付する（　　A　　）相当分の額を控除した額とする。

      \lawitem{一}
      （　　E　　）按分率
      
      \lawitem{二}
      積立金按分率

      \lawparagraph{2}
      前項の（　　D　　）は、当該年度における（　　F　　）等の総額に、当該年度において政府等が負担し、
      又は納付する（　　A　　）保険料相当分の合計額を加えた額とする。
    
    \lawparagraph{3〜5　（略）}
    \end{lawarticle}
  \end{lawstructure}  
\end{fillblank}

\begin{answerbox}
  \begin{enumerate}[A: ]
    \item 基礎年金拠出金
    \item 二分の一
    \item 共済各法
    \item 拠出金算定対象額
    \item 標準報酬
    \item 厚生年金保険給付費
  \end{enumerate}
\end{answerbox}

\begin{fillblank}[年金2 H28 問題1(4)]
  次は、被用者年金一元化の仕組みに関する記述である。以下の空欄を埋めなさい。

  \begin{itemize}
    \item 被用者年金一元化後の厚生年金は1・2階部分の保険料率および給付を統一した上で1・2 階部分の保険料収入および積立金を被用者全体の共通財源とした財政運営を行い、
    \begin{itemize}
      \item 共済組合等が徴収した厚生年金保険料および管理運用する1・2階積立金や標準報酬などの負担能力に応じて、共済組合等は年金特別会計の厚生年金勘定に（　　A　　）を納付し、
      \item 共済組合等が行う厚生年金の保険給付に要する費用等は同勘定から共済組合等に（　　B　　）として交付される

    \end{itemize}
    仕組みとなっている。

    \;

    \item 各実施機関から厚生年金勘定への（　　A　　）については徴収した厚生年金保険料および管理運用する1・2階積立金に応じて納付することを基本として次の2つの部分から構成されている。
    \begin{itemize}
      \item 標準報酬按分部分 \\
      厚生年金の標準報酬総額に対する実施機関ごとの標準報酬総額に応じた率にじたもの(標準報酬按分率)によって按分したもの。
      \item 積立金按分部分 \\
      厚生年金の積立金に対する実施機関ごとの1・2階相当積立金に応じた率に$( 1- （　　\text{C}　　） )$を乗じたものに（　　D　　）を乗じて得た率(積立金按分率)によって按分したもの。
      \item 支出費按分部分 \\
      厚生年金の給付費(共済年金の2階部分を含む)に （　　E　　） (国庫・公経済負担除く )を 加えたものに対する実施機関ごとの支出費に応じた率に（　　D　　）を乗じて得た率(支出費按分率)によって按分したもの
    \end{itemize}
    
  \end{itemize}
\end{fillblank}

\begin{answerbox}
  \begin{enumerate}[A: ]
    \item 拠出金
    \item 交付金
    \item 保険料財源比率
    \item 50/100
    \item 特別国庫負担額
  \end{enumerate}
\end{answerbox}

\newpage